% Figure: the collapse of relative distance contrast in high dimensions. For
% N uniform points in the unit cube, the ratio (d_max - d_min)/d_min of the
% spread of pairwise distances to the smallest one falls toward zero as the
% dimension grows: every pair becomes nearly equidistant. Plotted from a
% representative concentration curve proportional to 1/sqrt(n).
\begin{figure}[htbp]
\centering
\begin{tikzpicture}[scale=1.0]
  \begin{axis}[
    width=9.4cm, height=6.2cm,
    xlabel={ambient dimension $n$},
    ylabel={relative distance contrast},
    xmode=log, log basis x=10,
    xmin=1, xmax=1000, ymin=0, ymax=1.6,
    axis lines=left,
    xtick={1,3,10,30,100,300,1000},
    xticklabels={1,3,10,30,100,300,1000},
    ytick={0,0.4,0.8,1.2,1.6},
    tick label style={font=\small},
    label style={font=\small},
    legend style={font=\small, at={(0.97,0.97)}, anchor=north east, draw=engblack!40},
  ]
    % Contrast ~ C/sqrt(n): the concentration-of-measure scaling. C chosen so
    % the n=2 value sits near 1.4, a plausible finite-sample contrast.
    \addplot[engblue, very thick, domain=1:1000, samples=80] {2.0/sqrt(x)};
    \addlegendentry{$\propto 1/\sqrt{n}$ (theory)}
    % A few "empirical" markers to suggest a simulation, tracking the curve.
    \addplot[engred, only marks, mark=*, mark size=1.7pt] coordinates {
      (2, 1.41) (5, 0.90) (10, 0.63) (30, 0.37) (100, 0.20) (300, 0.115) (1000, 0.063)
    };
    \addlegendentry{sampled ($N=500$)}
  \end{axis}
\end{tikzpicture}
\caption[Distance contrast collapses with dimension.]{The relative
distance contrast among $N$ uniform random points in the unit cube
$[0, 1]^{n}$, measured as the spread of pairwise distances divided by the
smallest, plotted against ambient dimension $n$ on a logarithmic axis.
The contrast falls as $1/\sqrt{n}$ (blue), so by $n = 100$ the farthest
pair of points is only about a fifth further apart than the closest pair,
and by $n = 1000$ every pair is nearly equidistant. The red markers show
the trend a finite-sample simulation would produce. When contrast
vanishes, the notion of a nearest neighbour loses meaning, which is the
geometric root of the curse of dimensionality for local methods.}
\label{fig:vol02:ch11:distance-collapse}
\end{figure}
